\documentclass[11pt]{article}
\usepackage[T1]{fontenc}
\usepackage[utf8]{inputenc}
\usepackage{lmodern,microtype,geometry,amsmath,amssymb,amsthm,mathtools,booktabs}
\usepackage[colorlinks=true,linkcolor=blue,citecolor=blue,urlcolor=blue]{hyperref}
\geometry{margin=1in}
\hypersetup{pdftitle={Exponential mistakes and polynomial guarantees for circuit Occam learning},pdfauthor={ChatGPT-6 Astra; Claude; Kaarel Hanni; Dmitry Vaintrob}}
\newtheorem{theorem}{Theorem}
\newtheorem{lemma}[theorem]{Lemma}
\newtheorem{proposition}[theorem]{Proposition}
\newtheorem{corollary}[theorem]{Corollary}
\newtheorem{remark}[theorem]{Remark}
\newcommand{\E}{\mathbb E}
\newcommand{\Prb}{\mathbb P}
\title{Exponential mistakes and polynomial guarantees\\for circuit Occam learning}
\author{ChatGPT-6 Astra \and Claude (Anthropic) \and Kaarel H\"anni \and Dmitry Vaintrob}
\date{Research draft, 30 October 2026}
\begin{document}
\maketitle
\begin{abstract}
An online circuit Occam learner predicts using a smallest circuit
consistent with all previously revealed labels. Does a small target
circuit imply few mistakes when inputs arrive from changing or adversarial
regions? The answer depends decisively on how the learner selects among
equally small circuits.

For every deterministic selection rule based on the previous history,
the conjunction of $d$ variables admits a distinct-input ordering with at
least $2^d-d-2$ mistakes. We certify exact circuit minimality, including
unrestricted NAND circuits, where the target's minimum gate count is
$2d-2$. This gives an exponential lower bound in the target gate count.
The corresponding two-input-gate result was already proved in earlier
work; we present a short proof and a NAND certificate argument.

Conversely, freshly sampling uniformly among the smallest consistent
circuits gives $O(s^2\log(d+s))$ total expected mistakes for a target of
size $s$, on every nonanticipating adaptive input sequence. Majority
prediction within the minimum-size class has a deterministic bound of
the same order. For any history-based tie rule, arbitrary adaptive
interleaving of $K$ fixed input environments gives polynomial mistake
bounds when the environments have controlled disagreement probabilities;
in particular, arbitrary uniform subsets of the Boolean cube give
$O(K(s+1)\log(d+s+2)(d+1))$. Local random shuffling of the conjunction
example reduces exponential mistakes to quadratic expected mistakes.
We also prove bounds for continuing adaptive distribution drift and a
precise constant-fraction-elimination assumption. Exponentially small
distance from uniform, even with exponentially small total successive
drift, does not alone prevent exponentially many mistakes. All results concern
noiseless online feedback, not unlabeled transfer after training.
\end{abstract}

\section{The question and the prediction protocol}

Simplicity-based prediction can extrapolate because one compact rule may
apply in many regions of its input space. A worst-case online guarantee
asks for more: after every prediction the true label is revealed, and
the total number of mistakes should be polynomial in the complexity of
that rule, even if the queried region changes. We study this question for
exact minimum-circuit prediction, with unlimited computation.

Fix $d\ge1$ and a target $f:\{0,1\}^d\to\{0,1\}$. At round $t$, an
input $x_t$ is presented, the learner predicts, and $f(x_t)$ is revealed.
The label function never changes. Inputs may repeat, but exact consistency
makes a repeated input harmless. The conjunction-ladder lower bounds use
no repetitions; the near-uniform adaptive construction later uses repeats.
The mistake count is over the whole sequence, including an infinite
sequence where indicated. This is an online guarantee with feedback,
not a fixed-model test guarantee under an arbitrary new distribution.

A circuit is an acyclic list of gates with two initial constant wires
and the $d$ variable wires. Each gate reads earlier wires; fan-out and
sharing are unrestricted. The output is the final gate, and size is the
number of gates. We use either all sixteen two-input Boolean gates or
only two-input NAND gates, stating the basis explicitly. Binary AND/OR
and unary NOT are also covered by the first lower-bound proof.
No free inversions are allowed in the NAND-only result.
The target size $s=C(f)$ is the minimum gate count in the chosen basis.
All logarithms are natural unless a base is displayed. Every confidence
parameter $\delta$ below lies in $(0,1)$.

Three prediction rules must be distinguished:
\begin{enumerate}
\item A \emph{deterministic history-based selector} chooses one smallest
consistent circuit from previous observations and evaluates that circuit
on the next input. This includes a fixed ordering of circuit descriptions
to resolve ties. The circuit choice does not use the current input.
\item A \emph{fresh uniform selector} draws independently and uniformly
from the currently smallest consistent circuit descriptions each round.
The current input is fixed before that fresh draw, which is private until
then. An adaptive input chooser may use all previous predictions.
\item A \emph{minimum-class majority predictor} takes the majority label
of the smallest consistent descriptions on the current input. It may
equivalently select a smallest circuit giving that label, but its circuit
selection then uses the current input.
\end{enumerate}
Uniformity over descriptions need not mean uniformity over functions.
Our general counting bounds allow duplicate descriptions. None of these
rules is automatically the posterior of a Gaussian neural-network prior.

\paragraph{The main distinction.}
The first rule has an exponential worst-case lower bound, even for an
AND target. The second and third have polynomial upper bounds for every
allowed input order. Thus ``choose a smallest consistent circuit'' alone
does not specify a learner well enough to determine its online behavior.
The exponential construction can use a simply described order, so low
description complexity of the input sequence alone is not a sufficient
positive assumption. Our positive results for the first rule instead
give explicit stochastic conditions on how new inputs are selected.

\paragraph{Circuit counts.}
There are exactly
\[
 N_j=\prod_{i=1}^j16(d+i+1)^2
\]
size-$j$ descriptions in the all-binary-gate model, and the same product
without the factors $16$ for NAND. In particular, with $N_0=0$,
\begin{align}
 h_s:=\log\sum_{j=1}^sN_j
 &\le\log s+s\log16+2s\log(d+s+1),\label{eq:class-count}\\
 B_s:=s+\sum_{j=1}^s\log N_j
 &\le s+s(s+1)\log(4(d+s+1)).\label{eq:tier-count}
\end{align}
The learner need not know $s$; these are proof budgets. The realizing
target ensures that no consistent minimum-size choice ever exceeds $s$.

\paragraph{Relation to earlier work.}
The conjunction ladder appeared in the Vaintrob--H\"anni slides
\cite{slides}. H\"anni and Claude \cite{soi} proved its exact
minimum-circuit version for all binary gates and distinguished it from
Bayesian averaging. We reproduce that result with an elementary support
proof, then give a separate NAND-only proof. The positive arguments are
self-contained finite-class and version-space calculations, including
the standard halving principle \cite{littlestone}. They make the input,
tie-breaking, and complexity conventions explicit. These are conventional
proofs with independent AI-agent checks, not Lean formalizations.

% Includable section. Requires amsmath/amssymb/amsthm and theorem,
% lemma, proposition, and remark environments in the main document.
% The conjunction ladder is existing work, not a new result.
\section{An exponential lower bound for exact circuit minimization}
\label{sec:exact-occam-lower}

The conjunction ladder below is the construction from Appendix~K of
the Vaintrob--H\"anni slides \cite{slides}, proved for unrestricted
binary-gate circuits by H\"anni and Claude \cite{soi}.
We give an elementary support-based proof and make the prediction
protocol explicit. Its lower bound is exponential in the target's
\emph{minimum gate count}, and every hypothesis used by the learner
really is a globally smallest consistent circuit.

Throughout this section a circuit is a finite, acyclic, single-output
circuit over either all sixteen two-input Boolean gates or the standard
basis consisting of binary AND and OR and unary NOT. Constants are
available as source wires, and arbitrary fan-out is allowed. The output
must be a gate, so a circuit has at least one gate. Size counts every
gate, including NOT, and $C(f)$ denotes the minimum size of a circuit
computing $f$ on the full Boolean cube.

A \emph{history-only minimum-circuit learner} chooses a circuit $h_t$
of minimum gate count among all circuits consistent with the first
$t-1$ labeled examples, and predicts $h_t(x_t)$. Its choice of $h_t$
can depend on the entire past, but not on the current query $x_t$.
For deterministic learners this includes every fixed ordering of
equal-size circuit descriptions or equal-complexity functions.

\begin{lemma}[Support forces exact minimality]
\label{lem:lower-support}
A circuit of $s$ gates depends on at most $s+1$ input variables.
Suppose $h(1^d)=1$ and $h(x)=0$ whenever $|x|<k$.
Then $h$ depends on at least $k$ variables. If it depends on exactly
$k$ variables, it is the conjunction of those variables.
\end{lemma}
\begin{proof}
Prune gates having no path to the output, leaving $t\le s$ gates.
Of their at most $2t$ input slots, at least $t-1$ reference other
remaining gates. Thus at most $t+1\le s+1$ slots can hold distinct
input variables. This argument permits unary gates and shared wires.

Let $S$ be the essential-variable set of $h$, and let $1_S$ denote
its indicator in $\{0,1\}^d$. If $|S|<k$, then
$h(1_S)=h(1^d)=1$, contradicting the lower-weight negative labels.
If $|S|=k$, extend any assignment on $S$ other than $1_S$ by zeros
outside $S$. The resulting input has weight below $k$, so its label
is zero. The all-ones assignment on $S$ has label one. Hence
$h(x)=\bigwedge_{j\in S}x_j$.
\end{proof}

\begin{theorem}[Conjunction ladder for exact minimum circuits]
\label{thm:exact-occam-ladder}
Let $d\ge3$ and $f(x)=\bigwedge_{j=1}^d x_j$.
Then $C(f)=d-1$. For every deterministic history-only
minimum-circuit learner, there is a permutation of $\{0,1\}^d$
on which its total number of mistakes is at least
\[
  \sum_{k=2}^{d-1}\binom dk
  =2^d-d-2
  =2^{C(f)+1}-C(f)-3.
\]
The permutation may depend on the learner's tie rule.
\end{theorem}
\begin{proof}
First reveal $1^d$ with label one, then $0^d$ and all singleton
inputs with label zero. For $k=2,\ldots,d-1$, process the inputs of
weight $k$, choosing the order within the level as follows.

At the beginning of this phase, every input of weight less than $k$
has been revealed with label zero. At an intermediate step, let
$\mathcal R$ be the collection of $k$-sets whose indicator inputs
have not yet been queried. If $\mathcal R\ne\varnothing$, each
\[
  h_S(x)=\bigwedge_{j\in S}x_j,\qquad S\in\mathcal R,
\]
is consistent and has a circuit of $k-1$ gates. Every consistent
circuit has at least $k-1$ gates by Lemma~\ref{lem:lower-support}.
Consequently the \emph{global} minimum over all consistent circuits
is exactly $k-1$.

Every minimizing circuit depends on at most $k$ variables and at
least $k$ variables. The same lemma therefore says that it computes
$h_S$ for a $k$-set $S$. Consistency forces $S\in\mathcal R$.
Query $1_S$. The learner predicts one and the target label is zero.
Moreover, for another $k$-set $T$,
\[
  h_T(1_S)=1\quad\Longleftrightarrow\quad T\subseteq S
  \quad\Longleftrightarrow\quad T=S.
\]
Thus this query eliminates precisely one of the surviving
$k$-conjunctions. Repeat until the level is exhausted. Every input
in every level $k=2,\ldots,d-1$ is a mistake.

The query can be determined before prediction because the learner's
current hypothesis is a deterministic function of its past.
Equivalently, simulate the construction in advance to obtain a
fixed permutation for this learner. Finally, $f$ depends on all
$d$ variables, so Lemma~\ref{lem:lower-support} gives
$C(f)\ge d-1$; the usual AND tree attains equality.
\end{proof}

\paragraph{A fixed tie rule and a fixed input order.}
There is also a direct nonadaptive example. At each size $s\ge1$,
rank canonical left-associated conjunction circuits on $(s+1)$-sets
first, ordered lexicographically by their supports, and rank all
other circuits afterward in any fixed order. Choose the first
consistent circuit of minimum size. Present $1^d$, $0^d$, all
singletons, and then the levels of weight $2,\ldots,d-1$ in that
same lexicographic order. At every step after the singleton level,
the current hypothesis is the conjunction indexed by the next
input. The number of mistakes is exactly $2^d-d-2$. Both orderings
are fixed independently of the run. The theorem is stronger in a
different direction: every deterministic history-only tie rule
has some bad ordering, without asserting one ordering bad for
all tie rules.

\paragraph{Why the singleton level is omitted.}
Under the mandatory-output-gate convention, a one-gate
$x_i\wedge x_j$ is consistent with the initial positive example,
the zero vector, and every singleton negative. Therefore the
first level need not incur any mistakes. If instead an input or
constant wire may itself be a zero-gate output, the proof starts
at $k=1$ and yields $2^d-2$ mistakes after the two initial examples,
with the same target complexity $d-1$. This stronger count uses
a different, explicitly stated circuit convention.

\paragraph{The complexity measure matters.}
The theorem is literally exponential in the minimum number of
gates. For the explicit gate-list language that encodes size in
unary and then gate types and wire addresses, the target has
description length $b=\Theta(d\log d)$. In that measure the
conclusion is $2^{\Omega(b/\log b)}$, not a proved
$2^{\Omega(b)}$ bound. A constant-factor compilation into a
NAND-only basis does not automatically transfer the theorem:
compilation need not preserve which consistent circuits are
smallest. The next section proves the NAND-only statement directly.

\paragraph{Scope of the tie-breaking quantifier.}
The deterministic theorem does not permit an adversary to inspect
a hidden fresh random circuit before selecting its query. Nor does
it cover a deterministic learner that uses the current query to
choose between minimum circuits. For example, after receiving $x_t$,
a learner may take the majority label among all current minimum
circuits and then select a minimum circuit giving that label.
This still selects a smallest consistent circuit at every round.
Each mistake halves the candidates at the current size, or exhausts
that size. If $N_s$ bounds the number of size-$s$ descriptions, its
mistakes are at most
\[
  \sum_{s=1}^{C(f)}\bigl(1+\log_2 N_s\bigr)
  =O\!\left(C(f)^2\log(d+C(f))\right).
\]
Thus specifying only ``a smallest circuit each round'' leaves a
material ambiguity: the exponential lower bound concerns a single
deterministic hypothesis chosen from the past, before using the
current query.

\section{The same separation for pure NAND circuits}

The separation also holds when every gate is a two-input NAND. Size
means the number of NAND gates; intermediate results and input wires may
be reused without restriction. Constants may be allowed or omitted.
Negations and complemented edges are not free operations in this model.

\begin{lemma}[A positive certificate for a NAND circuit]
\label{lem:occam:nand-certificate}
Let a NAND circuit with $s\ge1$ gates compute $h$, and suppose
$h(1^d)=1$. There is a set $R\subseteq[d]$ with
$|R|\le\lfloor s/2\rfloor+1$ such that $h(x)=1$ whenever
$x_i=1$ for every $i\in R$. In particular, $h(1_R)=1$.
\end{lemma}

\begin{proof}
Evaluate every gate at $1^d$, calling it true or false according to its
value. Starting at the true output gate, retain one false-valued input
of each selected true NAND gate, and both true-valued inputs of each
selected false NAND gate. Recurse to the selected sources, stopping at
variables and constants. Identify repeated occurrences of the same
source or gate, but keep repeated edges when both input slots select
the same source. The resulting proof graph is a connected rooted
directed acyclic multigraph.

Let $T,F,L$ be its numbers of true gates, false gates, and distinct
source leaves, respectively. Every selected false gate has a selected
true parent, whereas every true gate selects only one input. Thus
$F\le T$: choose one true parent of each false gate; distinct false
gates cannot choose the same parent. There are $T+2F$ selected edges
and $T+F+L$ vertices. Connectedness gives
\[
 T+F+L-1\le T+2F,\qquad L\le F+1.
\]
Since $T+F\le s$ and $F\le T$, we obtain
$L\le\lfloor s/2\rfloor+1$.

Let $R$ consist of the distinct variable leaves, so $|R|\le L$.
Every one of these variables is one at $1^d$. Fixing them to one
forces every selected gate to retain its value at $1^d$, by induction
in circuit order: a selected false NAND has both inputs forced true,
and a selected true NAND has one input forced false. Constants retain
their values without any restriction on the input. The output
therefore remains one. The argument explicitly permits sharing and
repeated input wires.
\end{proof}

\begin{corollary}[Exact NAND complexity of conjunctions]
\label{cor:occam:nand-and}
For $k\ge2$, the conjunction of $k$ distinct variables has minimum
NAND-circuit size $2k-2$.
\end{corollary}

\begin{proof}
An AND of two signals is a NAND followed by a NAND with its two inputs
tied together. Replacing every gate in an AND tree gives $2k-2$ gates.
Conversely, every positive certificate for a conjunction must include
all its $k$ variables. The lemma gives
$k\le\lfloor s/2\rfloor+1$, hence $s\ge2k-2$.
\end{proof}

\begin{theorem}[Exponential mistakes for smallest NAND circuits]
\label{thm:occam:nand-ladder}
Fix $d\ge3$ and let $f(x)=\bigwedge_{i=1}^d x_i$. Consider any
deterministic learner which, before receiving the next input, chooses
a minimum-size NAND circuit consistent with all preceding examples.
There is a permutation of $\{0,1\}^d$, labeled throughout by $f$,
on which the learner makes at least
\[
 2^d-d-2
   =2^{s/2+1}-s/2-3
 \quad\hbox{mistakes},\qquad s=C_{\mathrm{NAND}}(f)=2d-2.
\]
The permutation may depend on the deterministic learner, but can be
fixed before the sequence is presented.
\end{theorem}

\begin{proof}
First present $1^d$ with label one, then every input of weight zero
or one with label zero. For $k=2,\ldots,d-1$, present the weight-$k$
inputs as follows.

At any point in stage $k$, every input of weight below $k$ has already
been labeled zero. If a consistent circuit has $r$ gates, the lemma
provides a positive certificate $R$ with
$|R|\le\lfloor r/2\rfloor+1$. Consistency requires $|R|\ge k$,
since otherwise $h(1_R)=1$ contradicts a previous label. Therefore
every consistent circuit has at least $2k-2$ gates.

As long as some weight-$k$ input $1_S$ remains unshown, the conjunction
$\bigwedge_{i\in S}x_i$ is consistent: it is one at $1^d$, zero below
weight $k$, and among weight-$k$ inputs it is one only at $1_S$.
It has $2k-2$ NAND gates. Thus the learner's minimum consistent size
is exactly $2k-2$. A positive certificate $R$ for its chosen circuit
has $|R|\le k$, and consistency forces $|R|\ge k$. Present $1_R$.
This input has not been shown before, the learner predicts one, and
the target label is zero because $|R|=k<d$.

Repeating this step presents every weight-$k$ input as a mistake.
Summing over $k=2,\ldots,d-1$ gives
$\sum_{k=2}^{d-1}\binom dk=2^d-d-2$. The initial points and these
levels comprise the whole hypercube. Because the learner is
deterministic, the constructed sequence can be simulated and fixed
in advance. The target labels never depend on the learner.
\end{proof}

\paragraph{A simple fixed tie rule and input order.}
For every $k\ge2$, fix the usual left-associated AND chain on a sorted
list of $k$ variable indices and replace each AND by its two-NAND
implementation. At size $2k-2$, place these canonical circuits first,
ordered lexicographically by their supports, and place all remaining
circuit descriptions afterward in any fixed order. At other sizes,
use any fixed order. Present $1^d$, then $0^d$, then the remaining inputs in
increasing Hamming weight, using lexicographic support order within
each level. At every stage $k\ge2$, the minimum consistent size is
$2k-2$, and the first surviving canonical conjunction has exactly the
support of the next input. This fixed tie rule and fixed input order
therefore realize the lower bound. No classification of all
minimum-size NAND hypotheses as conjunctions is required.

\paragraph{Conventions.}
The argument starts at $k=2$, so it applies both when outputs must be
gates and when direct wires or constants are allowed as zero-gate
outputs. If desired, the certificate lemma extends to those
zero-gate cases: a projection has a one-variable certificate and
constant one has the empty certificate. The result concerns gate
count. In a literal gate-list code with input addresses, the target's
description length is $\Theta(d\log d)$, so the lower bound becomes
$2^{\Omega(\ell/\log\ell)}$ in that description length.

% Requires amsmath, amssymb, amsthm and theorem/lemma/remark environments.
\section{Polynomial mistakes for fresh uniform ties and minimum-class majority}

Let $\mathcal C_j$ be a finite set of circuit descriptions of size $j$,
and write $N_j=|\mathcal C_j|$. Distinct descriptions computing the same
function are counted separately. Let $f$ have a description of size $s$.
All labels are noiseless, every label is revealed after prediction, and
all subsequent predictions must be consistent with every revealed label.
No algorithmic efficiency is asserted.

\subsection{Fresh uniform tie-breaking: arbitrary adaptive input order}

\begin{theorem}[Fresh uniform minimum-description prediction]\label{thm:uniform-min}
On each round, choose a fresh uniformly random description from the
currently consistent descriptions of minimum size, and predict with it.
The input may be chosen adaptively using the entire past and the learner's
previous predictions, but must be chosen before the fresh random choice
of the current description. Put
\[
 B_s=s+\sum_{j=0}^s\log\max\{1,N_j\}.
\]
Then the total number $M_\infty$ of mistakes over an arbitrarily long,
possibly infinite, input sequence satisfies
\[
 \mathbb E M_\infty\le B_s,
 \qquad
 \mathbb P\{M_\infty\ge(e-1)B_s+\log(1/\delta)\}\le\delta.
\]
For the all-binary-gate convention of this paper, and hence also for NAND,
\[
 B_s\le s+s(s+1)\log(4(d+s+1)),
\]
as in \eqref{eq:tier-count}. The empty size-zero class is omitted.
In particular this is a polynomial total-mistake guarantee for exact
minimum-size prediction under fully adaptive nonanticipating input order.
The learner need not know $s$.
\end{theorem}

\begin{proof}
Immediately before a label is revealed, let $j$ be the minimum surviving
size, let $n$ be the number of surviving descriptions at that size, and
let $q$ be the fraction of these descriptions wrong on the current input.
The conditional probability of the learner's mistake is exactly $q$.
If $q<1$, the surviving count at this same size becomes $n(1-q)$ and
\[
 q\le-\log(1-q)=\log n-\log(n(1-q)).
\]
Thus the sum of all such $q$ while size $j$ is the minimum is at most
$\log N_j$. If $q=1$, the size-$j$ class is exhausted and this transition
costs one. Minimum size never decreases. It never exceeds $s$, and the
size-$s$ class cannot be exhausted because the target survives. There are
therefore at most $s$ exhausted tiers. On every realization,
\[
 \sum_{t\ge1}q_t\le B_s.
\]
For the actual mistake indicators $Z_t$, conditional on the past and the
current input, $\mathbb E Z_t=q_t$ and
\[
 \mathbb E\exp\{Z_t-(e-1)q_t\}
 =e^{-(e-1)q_t}(1+(e-1)q_t)\le1.
\]
Consequently $\exp\{M_T-(e-1)\sum_{t\le T}q_t\}$ is a nonnegative
supermartingale, also when the next input is chosen from the full past.
Monotone convergence proves the expectation bound, and Ville's inequality
proves the time-uniform tail bound. The counting estimate is \eqref{eq:tier-count}.
\end{proof}

\begin{remark}[What the tie rule does and does not say]
Uniformity is over the specified finite set of minimum-size descriptions,
not automatically over distinct Boolean functions. Uniform sampling over
the latter also works if the corresponding finite counts are used.
The random choice must be fresh and private until the current input is
fixed. A single random ordering of descriptions, reused forever, is not
conditionally uniform after an adaptive input chooser has observed
predictions, so the above proof does not apply to that rule.
\end{remark}

\begin{corollary}[Minimum-tier majority]\label{cor:minimum-majority}
Predicting the majority label among the currently consistent descriptions
of minimum size, with arbitrary ties, makes at most $2B_s$ mistakes on
every input sequence. This can also be implemented as a query-dependent
choice of a minimum-size consistent circuit agreeing with that majority.
\end{corollary}

\begin{proof}
The same pathwise bound on $\sum q_t$ holds without any random prediction
rule. Each majority mistake has $q_t\ge1/2$.
\end{proof}


\subsection{The conjunction ladder under fresh uniform ties}
\begin{proposition}[Fresh uniform ties on the same ladder]
\label{prop:uniform-ties-ladder}
Use the all-binary-gate or AND/OR/NOT basis, with all gates charged.
Suppose that before every prediction the learner freshly samples
uniformly from all consistent circuits of minimum gate count.
Consider the level-by-level input order from
Section~\ref{sec:exact-occam-lower}, with the order
within each level chosen before the current fresh random draw.
For each $k=2,\ldots,d-1$, its expected mistakes in level $k$ are
exactly
\[
  H_{\binom dk}=\sum_{r=1}^{\binom dk}\frac1r.
\]
Thus its expected mistakes in these levels are
\[
  \sum_{k=2}^{d-1}H_{\binom dk}=\Theta(d^2),
\]
and the initial $d+2$ queries contribute at most $d+2$ further
mistakes.
\end{proposition}
\begin{proof}
Lemma~\ref{lem:lower-support} shows that the minimum circuits
in level $k$ compute precisely the surviving $k$-conjunctions.
Permuting input coordinates is a bijection between the size-$(k-1)$
circuits computing any two such conjunctions. They therefore have
equal numbers of minimum-size descriptions, in each of these two
bases. Since all descriptions computing a given function agree
on the entire observed history, uniformity over minimum circuit
descriptions induces uniformity over the surviving conjunctions.

If $r$ remain, any unqueried input of weight $k$ is predicted
incorrectly by exactly one of them. The conditional mistake
probability is $1/r$, and the revealed negative label removes
exactly that conjunction. Summing gives the harmonic number.
For the asymptotic estimate, $H_n\le1+\ln n$ gives an upper bound
$O(d^2)$. For $\Theta(d)$ central values of $k$,
$\binom dk=\exp(\Omega(d))$, so
$H_{\binom dk}\ge\ln\binom dk=\Omega(d)$, proving the lower bound.
\end{proof}

\section{Arbitrary history-dependent tie-breaking in changing environments}

The next results apply to every finite class $H$ containing $f$, with
$h=\log|H|$. On each round the learner selects a consistent member of $H$
using the previous history, before seeing the next sampled input. It may
randomize and its tie rule may depend on all previous source identities,
inputs, labels, and predictions. Minimum-size consistent circuits satisfy
this requirement with
$H=\bigcup_{j\le s}\mathcal C_j$, provided the choice of a minimum circuit
does not depend on the current sampled input. In this case
\[
 h\le\log s+s\log16+2s\log(d+s+1).
\]
The input-independent choice condition matters: a different consistent
hypothesis chosen after seeing each input can have much greater aggregate
error than any individual consistent hypothesis.

\begin{theorem}[Adaptively interleaved random-order environments]\label{thm:random-pools}
Fix $K$ finite input pools of sizes $m_1,\ldots,m_K\ge1$; repeated inputs
inside or between pools are allowed. Independently uniformly permute each
pool. At each round a scheduler, using the full observed past but no
unrevealed entries, chooses a nonexhausted pool and reveals its next
entry. It may stop at any time. Then every learner just specified satisfies
\[
 \mathbb E M\le K+2(h+1)\sum_{i=1}^K H_{m_i-1},
 \qquad H_0=0.
\]
For $0<\delta<1$, put
\[
 a=h+\log\frac{2\sum_i m_i}{\delta},\qquad
 B_\delta=K+2a\sum_i H_{m_i-1}.
\]
With probability at least $1-\delta$,
\[
 M\le(e-1)B_\delta+\log(2/\delta).
\]
In particular, if each pool has at most $2^d$ entries, expectation is
$O(K(h+1)(d+1))$, and high probability adds only polynomial dependence on
$\log(K/\delta)$ and $d$.
\end{theorem}

\begin{proof}
For pool $i$, let $R_{i,n}$ be the maximum pool error fraction of a member
of $H$ agreeing with the first $n$ entries in that pool, ignoring all other
labels. A candidate with error fraction $r$ survives that prefix with
probability at most $(1-r)^n\le e^{-nr}$. Thus
\[
 \mathbb P(R_{i,n}>u)\le\min\{1,e^{h-nu}\},
 \qquad \mathbb E R_{i,n}\le(h+1)/n.
\]
At the next visit to pool $i$, the selected predictor is consistent with
that prefix; all of its pool errors lie in the $m_i-n$ remaining entries.
Its conditional error probability is at most
$m_iR_{i,n}/(m_i-n)$. Charging each used local count $(i,n)$ once and
enlarging to all possible local counts gives the expectation bound because
\[
 \sum_{n=1}^{m_i-1}\frac{m_i}{n(m_i-n)}=2H_{m_i-1}.
\]
The first visit to each pool costs at most one. This argument is unchanged
by arbitrary adaptive interleaving because each unused next entry remains
uniformly distributed among its own pool's remaining entries.

For high probability, a union bound makes all inequalities
$R_{i,n}\le a/n$, $1\le n<m_i$, hold except with probability $\delta/2$.
Stop the mistake process immediately before using a failed prefix. Until
that stopping time its total conditional error probabilities are at most
$B_\delta$. The exponential supermartingale used above bounds its mistakes
by $(e-1)B_\delta+\log(2/\delta)$ except with probability $\delta/2$.
On the simultaneous good-prefix event the process never stops, proving
the claim.
\end{proof}

\begin{theorem}[Interleaved iid environments]
Fix $K$ distributions $D_1,\ldots,D_K$. A scheduler uses the full past to
choose a source, then receives a fresh independent input from that source.
For every source having at least one nonzero disagreement risk in $H$, let
\[
 \varepsilon_i=
 \min\{D_i(g\ne f):g\in H,\ D_i(g\ne f)>0\}.
\]
Omit and reindex sources on which all hypotheses agree with $f$; $K$
now counts only the remaining sources. If there are none, $M_\infty=0$
and no probability budget is needed. Otherwise,
\[
 \mathbb E M_\infty\le
 \sum_{i=1}^K\{1+(h+1)(1+\log(1/\varepsilon_i))\}.
\]
For a high-probability bound set
\[
 J_i=1+\lceil\log_2(1/\varepsilon_i)\rceil,\quad
 L=\sum_iJ_i,\quad a=h+\log(2L/\delta),\quad
 B_\delta=(L+K)a+2K.
\]
Then, with probability at least $1-\delta$,
\[
 M_\infty\le(e-1)B_\delta+\log(2/\delta).
\]
On a finite support with smallest positive atom $p_i$, one may replace
$\varepsilon_i$ by $p_i$. In particular, sources uniform on arbitrary
subsets of $\{0,1\}^d$ give an expectation of $O(K(h+1)(d+1))$ even with
unlimited adaptive switching among disjoint regions.
\end{theorem}

The source need not be observed by the learner. In particular the theorem
covers arbitrary predictable mixture weights over the fixed $D_i$: choose
a latent source from the current weights and then draw from that source.
There is no bound on the number of switches or on the distance between
the source distributions.

\begin{proof}
Couple source $i$ to its own iid stream and let $R_{i,n}$ be the maximum
$D_i$-risk of a hypothesis surviving the first $n$ samples of that stream.
The same union bound gives
$\mathbb P(R_{i,n}>u)\le\min\{1,e^{h-nu}\}$.
The variable $R_{i,n}$ is either zero or at least $\varepsilon_i$.
For $u>0$, monotonicity of the integrand gives
\[
 \sum_{n\ge1}\min\{1,e^{h-nu}\}
 \le\int_0^\infty\min\{1,e^{h-tu}\}\,dt=(h+1)/u.
\]
Consequently Tonelli's theorem, using $\varepsilon_i$ in place of $u$
for $0<u<\varepsilon_i$, yields
\[
 \sum_{n\ge1}\mathbb E R_{i,n}
 \le(h+1)\left\{\int_0^{\varepsilon_i}\frac{du}{\varepsilon_i}
                   +\int_{\varepsilon_i}^1\frac{du}{u}\right\}
 =(h+1)(1+\log(1/\varepsilon_i)).
\]
The initial local round has risk at most one. On every visit to source $i$
after $n$ samples, the current consistent predictor has risk at most
$R_{i,n}$, regardless of observations from other sources. Each local count
is used at most once. Summing these bounds proves the expectation claim.

For high probability, for each $1\le j\le J_i$ put
$n_{i,j}=\lceil a2^j\rceil$. A union bound gives, simultaneously,
\[
 R_{i,n_{i,j}}\le2^{-j}
\]
except with probability $Le^{h-a}=\delta/2$. Since
$2^{-J_i}<\varepsilon_i$, the last inequality implies
$R_{i,n_{i,J_i}}=0$. By monotonicity of $R_{i,n}$, the initial block costs
at most $2a+1$ and each subsequent block costs at most $a+2^{-j}$, so
\[
 \sum_{n\ge0}R_{i,n}\le(J_i+1)a+2.
\]
Stop immediately before using a prefix which violates the applicable
block bound. Until stopping, the conditional error sum is at most
$B_\delta$. The same exponential-supermartingale bound and the good-event
union bound establish the stated infinite-horizon probability guarantee.
\end{proof}

\begin{remark}[Bounded arbitrary departures]
Either environment theorem remains valid with at most $A$ extra arbitrary
input rounds, with an additive $A$ in the actual mistake bound and its
expectation. The extra noiseless observations only shrink the version
space. The clean source counters and their latent orders are unchanged.
This includes a finite adversarial excursion followed by arbitrarily
many further source switches.
\end{remark}


\subsection{Randomizing within Hamming layers restores a polynomial bound}

For circuits over all two-input Boolean gates, or over the
AND/OR/NOT basis, a simple randomization of the input order changes
the same example. Let $f(x)=\bigwedge_{i=1}^d x_i$, present $1^d$
and $0^d$ first, and then present the remaining Hamming levels in
increasing order. Independently for each level, use a uniformly
random permutation of that level. The learner remains deterministic
and chooses a smallest consistent circuit using any rule depending
only on the preceding history.

For $k\ge2$, after all smaller levels have been labeled zero, the
globally minimum consistent functions are precisely the surviving
$k$-variable conjunctions. Here is the elementary reason. If $h$
depends on the variables in $R$, then $h(1_R)=h(1^d)=1$, so
$|R|\ge k$. A circuit with $r$ binary gates depends on at most
$r+1$ variables. Surviving $k$-conjunctions attain size $k-1$;
at that minimum size, $h$ depends on exactly $k$ variables.
Every non-all-ones assignment on these variables has a zero
extension of weight below $k$, so $h$ is their conjunction.
The same argument applies to AND/OR/NOT circuits with all gates
charged, since unary gates add no dependencies.

Suppose that $m$ weight-$k$ inputs remain. The learner's chosen
conjunction is wrong at precisely one of them, its own support
indicator. Conditional on the complete preceding history, the next
input is uniform over the $m$ remaining inputs. Its conditional
mistake probability is therefore exactly $1/m$, regardless of the
history-dependent deterministic tie rule. Consequently stage $k$
has exactly
\[
 H_{\binom dk}=\sum_{m=1}^{\binom dk}\frac1m
\]
expected mistakes. The first two inputs and the weight-one level
contribute between zero and $d+2$ mistakes, so
\[
 \sum_{k=2}^{d-1} H_{\binom dk}
 \;\le\; \mathbb E M
 \;\le\; d+2+\sum_{k=2}^{d-1} H_{\binom dk}
 \;=\;\Theta(d^2).
\]
For the upper bound, $H_N\le1+\ln N\le1+d\ln2$.
For the lower bound, for each integer
$d/4\le k\le d/2$,
$\binom dk\ge(d/k)^k\ge2^{d/4}$, hence
$H_{\binom dk}\ge(d/4)\ln2$; there are $\Omega(d)$
such levels. The assertion is asymptotic in $d$.

Thus the order still progresses strictly from low to high Hamming
weight after the initial positive example; only the order inside
each level is randomized. This suffices to reduce the deterministic
learner's expected mistake count from exponential to quadratic in
the target's minimum gate count $d-1$. The exact phase expectations
hold for every deterministic history-only tie rule. No analogous
upper bound is asserted here for pure NAND circuits, because the
NAND proof uses certificates rather than a classification of every
minimum consistent function.

% Standalone includable section. Requires amsmath, amssymb and a theorem environment.
\section{Adaptive shifts with a persistent reference distribution}
\label{sec:reference-shifts}

A deterministic choice among minimum circuits also has a polynomial
all-time mistake guarantee under a useful class of non-i.i.d. processes.
The condition is that the next-input distribution may adapt to the whole
past, but remains comparable to one fixed reference distribution.  In
particular, it continues to sample every reference region often enough.
The reference need not be uniform.

Let $\mathcal H$ be a finite class of $N$ Boolean hypotheses containing a
fixed target $h_\star$.  Write $\mathcal F_t$ for the information available
after round $t$, and suppose the learner chooses an
$\mathcal F_{t-1}$-measurable $h_t\in\mathcal H$ that is consistent with
all preceding examples.  The choice may depend arbitrarily on the past,
including learner randomness included in the filtration.  Its circuit
must be chosen before seeing the current input.  Labels are always
$h_\star(X_t)$.  Denote the conditional law of $X_t$ given $\mathcal F_{t-1}$ by $P_t$.
The comparison below is with respect to this same filtration, including
any learner randomness revealed before the current input.

\begin{proposition}[Adaptive reference comparability]
\label{prop:adaptive-reference}
Suppose that, almost surely, for every round $t$ and measurable set $A$,
\[
 a\mu(A)\le P_t(A)\le b\mu(A),
 \qquad 0<a\le1\le b<\infty,
\]
where $\mu$ is a fixed probability distribution.  Suppose also that every
positive disagreement probability
$\mu\{h\ne h_\star\}$, for $h\in\mathcal H$, is at least $0<\gamma\le1$.
Then the total number $M_\infty$ of mistakes satisfies
\[
 \mathbb E M_\infty
 \le b\left[1+\frac{\log N+1}{a}
                    \left(1+\log\frac1\gamma\right)\right].
\]
In particular, $M_\infty<\infty$ almost surely.
\end{proposition}

\begin{proof}
For $h\in\mathcal H$, put $p(h)=\mu\{h\ne h_\star\}$ and let
$\mathcal V_n$ be the hypotheses consistent with the first $n$ examples.
The conditional lower bound implies
\[
 \Pr(h\in\mathcal V_n)\le(1-ap(h))^n\le e^{-anp(h)}.
\]
Set $R_n=\max_{h\in\mathcal V_n}p(h)$.  By a union bound,
\[
 \Pr(R_n>u)\le\min\{1,Ne^{-anu}\}, \qquad u>0.
\]
For $0<u<\gamma$, the same bound holds with $u$ replaced by $\gamma$,
because each positive $p(h)$ is at least $\gamma$.
For any $q>0$, monotonicity gives
\[
 \sum_{n=0}^{\infty}\min\{1,Ne^{-anq}\}
 \le 1+\int_0^\infty\min\{1,Ne^{-atq}\}\,dt
 =1+\frac{\log N+1}{aq}.
\]
Integrating the tail probability over $0<u<1$, splitting at $\gamma$,
and using Tonelli's theorem therefore yields
\[
 \mathbb E\sum_{n=0}^{\infty}R_n
 \le 1+\frac{\log N+1}{a}
                     \left(1+\log\frac1\gamma\right).
\]
Conditional on $\mathcal F_{t-1}$, the learner's mistake probability is
at most $b\,p(h_t)\le bR_{t-1}$.  Summation proves the claim.
\end{proof}

For the uniform reference on $\{0,1\}^d$, one may take
$\gamma=2^{-d}$.  If hypotheses have binary descriptions of at most $L$
bits, then $N<2^{L+1}$, giving
\[
 \mathbb E M_\infty
 =O\!\left(\frac ba(L+1)(d+1)\right).
\]
This applies to an exact minimum-description consistent circuit, since a
realizing target of description length $L$ keeps every selected circuit
within that class.  For minimum gate count $s$, standard circuit counting
gives $\log N=O((s+1)\log(d+s+2))$.  No efficient algorithm for finding the exact minimum
is asserted.  For a nonuniform finite reference, $\gamma$ may instead be
bounded below by its smallest positive atom; polynomial bounds hold when
$\log(1/\gamma)$ and $b/a$ are polynomial in the problem parameters.

For example, a process of the form
\[
 P_t=(1-\eta)\mu+\eta Q_t, \qquad 0\le\eta<1,
 \qquad Q_t(A)\le B\mu(A),
\]
may choose $Q_t$ adaptively.  It satisfies the proposition with
$a=1-\eta$ and $b=1-\eta+\eta B$.  It need never become i.i.d. or converge
to a limiting distribution.

\paragraph{Summably many departures from the comparison condition.}
The same conclusion has a robust extension.  Suppose there are
nonnegative $\mathcal F_{t-1}$-measurable $\varepsilon_t$ such that
\[
 a\mu(A)-\varepsilon_t\le P_t(A)
 \le b\mu(A)+\varepsilon_t,
 \qquad \sum_{t\ge1}\varepsilon_t\le V\quad\hbox{almost surely}.
\]
Then
\[
 \mathbb E M_\infty
 \le V+b\left[1+\frac{\log N+V+1}{a}
                         \left(1+\log\frac1\gamma\right)\right].
\]
Indeed, for each fixed $h$, the process
\[
 \mathbf1\{h\in\mathcal V_n\}
 \exp\!\left(anp(h)-\sum_{t=1}^n\varepsilon_t\right)
\]
is a nonnegative supermartingale.  Thus
$\Pr(h\in\mathcal V_n)\le e^{V-anp(h)}$.
The preceding proof applies with $N$ replaced by $Ne^V$, while the
conditional mistake probability is at most $bR_{t-1}+\varepsilon_t$.
In particular, with $a=b=1$ this covers
$\sum_t\|P_t-\mu\|_{\mathrm{TV}}\le V$.
Here total variation is $\sup_A|P(A)-Q(A)|$.
This assumption controls cumulative distance from a fixed reference; it
is different from bounding only the successive changes
$\sum_t\|P_{t+1}-P_t\|_{\mathrm{TV}}$.


% Includable exact-minimum-circuit obstruction for upper-only smoothness.
\subsection{Tiny persistent shifts can defeat deterministic minimum circuits}
\label{sec:smooth-lower}

An upper density bound alone does not suffice even for exact minimum
circuits.  The failure can persist when every conditional input law is
exponentially close to uniform and, for a fixed ranking of circuits, the
sum of successive distribution changes is exponentially small.
This construction uses repeated inputs.

\begin{proposition}[Upper smoothness does not ensure few mistakes]
\label{prop:upper-smooth-lower}
Use a fixed finite fan-in-two gate basis containing AND, and count gates.
Let $d\ge4$ be divisible by four, let $\mu$ be uniform on
$\{0,1\}^d$, and put $\delta=e^{-d/8}$.
For every deterministic, history-only exact-minimum consistent circuit
rule, there is an adaptive input process with target
$f(x)=\bigwedge_{j=1}^d x_j$, of minimum size $d-1$, such that
\[
 P_t(A)\le\frac{\mu(A)}{1-\delta}\quad\hbox{for every set }A,
 \qquad \|P_t-\mu\|_{\rm TV}\le\delta
 \quad\hbox{for every round }t,
\]
yet the learner makes at least $2^{d/2}$ mistakes almost surely.
If the rule selects the first consistent minimum circuit in a fixed
ranking, the same process also satisfies
\[
 \sum_{t\ge1}\|P_{t+1}-P_t\|_{\rm TV}
 \le\frac{\delta}{1-\delta}\quad\hbox{almost surely}.
\]
There is a deterministic horizon $T$ with $\log T=O(d)$ at which
$\mathbb E M_T\ge2^{d/2-1}$.
\end{proposition}

\begin{proof}
Put $k=3d/4$, let $R=\{x:|x|<k\}$, and write
$r=|R|$, $m=\binom dk$, and $N=2^d$.
First draw independently and uniformly from $B=R\cup\{1^d\}$ until
every point of $B$ has appeared.  This stopping time is finite almost
surely.  These examples force any consistent function to depend on at
least $k$ coordinates: a function depending on a set $S$ of fewer than
$k$ coordinates cannot distinguish $1^d$ from $1_S\in R$.
A circuit with $s$ binary gates depends on at most $s+1$ coordinates.
Every $k$-variable conjunction is consistent and has $k-1$ gates, so
minimum size is exactly $k-1$.

Moreover, every minimum circuit computes a $k$-variable conjunction.
It depends on exactly $k$ coordinates, say $S$, and is one on their
all-ones assignment.  Every other assignment to $S$ occurs on a point
of weight below $k$ and is labeled zero.  Its function is therefore
precisely $\bigwedge_{j\in S}x_j$.

While any $k$-conjunction remains consistent, let $S_t$ be the support of
the learner's current minimum circuit and choose $P_t$ uniform on
$B\cup\{1_{S_t}\}$.  The current circuit is determined by the past, so
this defines a valid adaptive input process.
A draw from $B$ causes no mistake and eliminates no remaining
$k$-conjunction.  A draw of $1_{S_t}$ has conditional probability
$1/(r+2)$, causes a mistake, and eliminates exactly that conjunction:
two distinct $k$-subsets are not contained in one another.  Therefore
minimum size remains $k-1$ until all $m$ conjunctions have been
eliminated.  Those $m$ errors occur almost surely.  Thereafter retain the
final distribution, whose entire support has already been observed.

Hoeffding's inequality gives
\[
 1-r/N=\Pr\{\operatorname{Bin}(d,1/2)\ge3d/4\}
 \le e^{-d/8}=\delta.
\]
Every conditional law is uniform on a set of size $r+1$ or $r+2$.
Its density relative to $\mu$ is at most $N/(r+1)\le(1-\delta)^{-1}$,
and its total variation distance from $\mu$ is at most
$1-(r+1)/N\le\delta$, including the initial and final stages.
Under a fixed ranking, repeated points of $B$ leave the selected circuit
unchanged.  There are $m$ hard phases.  The initial switch and the
$m-1$ switches between their laws each have total variation $1/(r+2)$;
there is no final switch.  Thus total successive variation is
$m/(r+2)\le\delta/(1-\delta)$, using $m\le N-r\le N\delta$.

Finally, $m=\binom{d}{d/4}\ge(d/(d/4))^{d/4}=2^{d/2}$.
This construction permits repeated inputs; the earlier exponential
lower bounds use a permutation with no repeats. The expected time through the initial stage and all hard phases is
\[
 (r+1)H_{r+1}+m(r+2),
\]
where $H_n$ is the $n$th harmonic number.  Taking $T$ to be the ceiling of
twice this expectation gives $\mathbb E M_T\ge m/2$ by Markov's
inequality, and $\log T=O(d)$.  This horizon is exponential in $d$;
only its logarithm is linear in $d$.
\end{proof}

\section{A precise genericity condition}

This final result applies to a particular simplest-consistent selector
with no prescribed tie-breaking, but its assumption is target-relative
and must be justified separately for the application. Let $V_t$ be all
surviving descriptions no longer than the target description, let $N$ be
their initial count, and let $q_t$ be the fraction of $V_t$ wrong on the
current input. The target survives, so $q_t<1$ and, pathwise,
\[
 \sum_t q_t\le\sum_t-\log(1-q_t)\le\log N.
\]
Let $Z_t$ indicate a mistake by the selected simplest description.

\begin{theorem}[Conditional genericity]\label{thm:genericity}
Suppose, for a constant $0<c\le1$, on every round with positive conditional
mistake probability,
\[
 \mathbb E[q_t\mid Z_t=1,\mathcal F_{t-1}]\ge c.
\]
Equivalently, without conditioning on null events, assume
$\mathbb E[Z_tq_t\mid\mathcal F_{t-1}]
 \ge c\mathbb E[Z_t\mid\mathcal F_{t-1}]$.
Then
\[
 \mathbb E M_\infty\le\frac{\log N}{c},\qquad
 \mathbb P\left\{M_\infty\ge
 \frac{\log N+\log(1/\delta)}{c(1-e^{-1})}\right\}\le\delta.
\]
It suffices that
$\mathbb P(q_t\ge\alpha\mid Z_t=1,\mathcal F_{t-1})\ge\beta$
with $c=\alpha\beta$. No independence, stationarity, or iid input
assumption is needed.
\end{theorem}

\begin{proof}
The expectation follows by summing
$c\mathbb E Z_t\le\mathbb E[Z_tq_t]\le\mathbb E q_t$.
For the tail let
$\theta=-\log(1-c(1-e^{-1}))\ge c(1-e^{-1})$.
Convexity on $[0,1]$ gives $e^{-q}\le1-(1-e^{-1})q$. Thus
\[
 \mathbb E[e^{\theta Z_t-q_t}\mid\mathcal F_{t-1}]
 \le1-p_t+e^\theta\{p_t-(1-e^{-1})\mathbb E[Z_tq_t\mid\mathcal F_{t-1}]\}
 \le1,
\]
where $p_t=\mathbb E[Z_t\mid\mathcal F_{t-1}]$.
Hence $\exp\{\theta M_T-\sum_{t\le T}q_t\}$ is a nonnegative
supermartingale. Since $\sum_tq_t\le\log N$, Ville's inequality proves
the result, uniformly over all finite horizons and therefore at infinity.
\end{proof}

\paragraph{Interpretation.}
These are three distinct positive routes. Fresh uniform minimum-tier
sampling supplies a theorem for arbitrary adaptive inputs by changing the
tie rule. Interleaved environments give an explicit structured input class
for every history-dependent tie rule, with a logarithmic domain-size
overhead. Conditional genericity directly formalizes the proposed
constant-fraction-elimination intuition. The environment theorems do not
assert that a constant fraction is removed on each mistake; they prove a
polynomial total by a different survival argument. None of these claims
by itself establishes arbitrary-order guarantees for a specified Gaussian
neural posterior.


\section{What follows, and what remains open}

The exact-minimality issue in the AND example is resolved for both
unrestricted binary-gate circuits and NAND-only circuits. No difficult
general circuit lower bound is needed: support size proves the former,
and a small positive certificate proves the latter. The target labels
are fixed in advance. For each deterministic history-based learner, the
adversarial construction can be simulated in advance into one fixed
permutation; it does not assert one permutation that defeats all learners.

There are rigorous positive results without changing that deterministic
selection rule. Independent shuffling inside each Hamming-weight level
already changes the conjunction example from exponential to quadratic
expected mistakes in the binary-gate basis. The environment theorems
allow unrelated and disjoint regions, arbitrary adaptive switches and
mixtures, and infinitely many rounds. Their polynomial bound charges the
number of fixed environments and their smallest positive disagreement
probabilities. Persistent-reference bounds permit continuing adaptive
drift instead. By contrast, tiny but persistent omissions of carefully
selected rare inputs can still produce exponentially many mistakes, even
with exponentially small distance from uniform and total successive drift.
These are explicit sufficient conditions and obstructions, not a measure
under which most conceivable input sequences are typical.

The genericity premise captures a different intuition: errors of the
chosen simplest rule should often also be errors of many surviving
simple rules. Its potential argument is exact, but verifying the premise
for realistic deterministic data streams remains a substantive task.
The randomized minimum-size theorem needs no such premise; it pays a
separate counting budget at each complexity level. Ordinary Bayesian
prediction under a fixed circuit-description prior has yet another
potential: cumulative log loss is bounded by the target's negative log
prior mass. These mechanisms should not be identified with each other.

Nothing here establishes an arbitrary-order polynomial mistake bound for
the particular Gaussian neural-network learners studied elsewhere.
Nor does the exponential circuit-Occam example disprove such a bound for
their posterior mixtures. Transferring a result requires controlling the
actual prediction law, not merely the complexity of one selected mode.

\begin{thebibliography}{9}
\bibitem{slides} Dmitry Vaintrob and Kaarel H\"anni.
\emph{Neural-network Bayesian generalization}, presentation, Appendix K.
\url{https://kaarelh.github.io/ai/nnbayes-generalization.pdf}.
\bibitem{soi} Kaarel H\"anni and Claude (Anthropic).
\emph{Solomonoff Operator Induction: Definitions, Versions and Basic Theorems}.
October 2026. Section ``Circuit SOI'' and its appendix proofs.
\url{https://github.com/kaarelh/AI-works/blob/main/solomonoff-operator-induction/paper/main.pdf}.
\bibitem{littlestone} Nick Littlestone.
Learning quickly when irrelevant attributes abound: A new linear-threshold
algorithm. \emph{Machine Learning} 2, 285--318, 1988.
\end{thebibliography}
\end{document}
